\documentclass{article}
\usepackage{amsmath}
\usepackage{booktabs} % よりきれいな表にするため
\usepackage{array}

\begin{document}

\begin{table}[h]
\centering
\setlength{\tabcolsep}{6pt}
\renewcommand{\arraystretch}{1.2}
\begin{tabular}{c|ccc|ccc}
\toprule
\textbf{現在の状態} & \multicolumn{3}{c|}{\textbf{次状態 } $q^{(t+1)}$} & \multicolumn{3}{c}{\textbf{出力 } $z$} \\
 & 00 & 01 & 10 & 00 & 01 & 10 \\
\midrule
$q_0$ & $q_0$ & $q_1$ & $q_5$ & 0 & 0 & 0 \\
$q_1$ & $q_1$ & $q_2$ & $q'_0$ & 0 & 0 & 1 \\
$q_2$ & $q_2$ & $q_3$ & $q'_1$ & 0 & 0 & 1 \\
$q_3$ & $q_3$ & $q_4$ & $q'_2$ & 0 & 0 & 1 \\
$q_4$ & $q_4$ & $q_5$ & $q'_3$ & 0 & 0 & 1 \\
$q_5$ & $q_5$ & $q'_0$ & $q'_4$ & 0 & 1 & 1 \\
$q'_0$ & $q_0$ & $q_1$ & $q_5$ & 0 & 0 & 0 \\
$q'_1$ & $q_1$ & $q_2$ & $q'_0$ & 0 & 0 & 1 \\
$q'_2$ & $q_2$ & $q_3$ & $q'_1$ & 0 & 0 & 1 \\
$q'_3$ & $q_3$ & $q_4$ & $q'_2$ & 0 & 0 & 1 \\
$q'_4$ & $q_4$ & $q_5$ & $q'_3$ & 0 & 0 & 1 \\
\bottomrule
\end{tabular}
\caption{状態遷移表}
\end{table}

\end{document}
